\section*{Capítulo \ref{ch:b1:wave-propagation} --- Propagação de sinais: ondas progressivas e estacionárias}

\begin{solution}{exo:b1:wave-propagation:1}
$d = 340 \times 6.0 = \qty{2.0}{km}$. Ida e volta: profundidade $= 1500 \times
0.40/2 = \qty{300}{m}$.
\end{solution}

\begin{solution}{exo:b1:wave-propagation:2}
$\lambda = 340/440 = \qty{0.77}{m}$; $\Delta\varphi = 2\pi \times
0.20/0.773 = \qty{1.6}{rad}$ ($93^\circ$); em \omterm{def:b1:wave-propagation:signal}{fase} para separações
$n\lambda = n \times \qty{0.77}{m}$.
\end{solution}

\begin{solution}{exo:b1:wave-propagation:3}
$c = \sqrt{80/4.0\times10^{-3}} = \qty{141}{m/s}$; $\lambda = c/f =
\qty{0.71}{m}$.
\end{solution}

\begin{solution}{exo:b1:wave-propagation:4}
$444 - 440 = 4$ batimentos por segundo; período da envoltória $1/4 = \qty{0.25}{s}$.
Afrouxar abaixa a corda: \qty{441}{Hz} (chegar a \qty{439}{Hz} teria
exigido passar pelo silêncio em \qty{440}{Hz}); verifique afrouxando
um pouco mais: os batimentos devem sumir e depois voltar.
\end{solution}

\begin{solution}{exo:b1:wave-propagation:5}
$f_1 = c/2L = 240/1.2 = \qty{200}{Hz}$; $400$, \qty{600}{Hz}. Terceiro modo:
nós em $0$, $0.20$, $0.40$, \qty{0.60}{m}. Toque a \qty{20}{cm} (ou
\qty{40}{cm}) de uma extremidade: um nó ali conserva apenas $n = 3, 6, \dots$, e
a corda soa em \qty{600}{Hz}.
\end{solution}

\begin{solution}{exo:b1:wave-propagation:6}
$\lambda = \qty{0.40}{m}$. No segmento, à distância $x$ do alto-falante
1, $\delta = (2.0 - x) - x = 2.0 - 2x$; silêncio para $\delta = \pm\lambda/2,
\pm 3\lambda/2, \dots$: $x = 0.1, 0.3, \dots, 1.9$ m, a cada \qty{20}{cm};
o ponto médio é forte. Na mediatriz $\delta = 0$: forte em todo o percurso (só
com o decaimento em $1/r$). Em oposição de \omterm{def:b1:wave-propagation:signal}{fase}: forte e fraco trocam --- silêncio no
ponto médio e a cada \qty{20}{cm} a partir dele.
\end{solution}

\begin{solution}{exo:b1:wave-propagation:7}
$i = \lambda D/a = 5.89\times10^{-7} \times 1.5/2.0\times10^{-4} =
\qty{4.4}{mm}$; azul: \qty{3.4}{mm}. Em $\pm\qty{1.0}{cm}$: franjas em
$0$, $\pm 4.4$, $\pm 8.8$ mm --- cinco.
\end{solution}

\begin{solution}{exo:b1:wave-propagation:8}
Iguais: $I_{\max} = 4I_0$, $I_{\min} = 0$, contraste $1$. $I_0$ e $I_0/4$:
$I_{\max} = I_0(1 + \tfrac14 + 2 \times \tfrac12) = 2.25I_0$, $I_{\min} =
I_0(1.25 - 1) = 0.25I_0$; contraste $2.0/2.5 = 0.8$.
\end{solution}

\begin{solution}{exo:b1:wave-propagation:9}
Fechado--aberto: $f = (2n + 1)c/4L = 170$, $510$, \qty{850}{Hz} (só
harmônicos ímpares: o clarinete). Aberto--aberto: $nc/2L = 340$, $680$,
\qty{1020}{Hz} (a flauta).
\end{solution}

\begin{solution}{exo:b1:wave-propagation:10}
$v_\varphi = \sqrt{9.81 \times 10/2\pi} = \qty{3.9}{m/s}$ e
\qty{12.5}{m/s} para \qty{100}{m}. A ondulação longa primeiro: $10^6/12.5 =
\qty{22}{h}$ contra $10^6/3.95 = \qty{70}{h}$, dois dias antes. A velocidade
depende do \omterm{prop:b1:wave-propagation:sinusoidal}{comprimento de onda}: um mar misturado se separa por \omterm{prop:b1:wave-propagation:sinusoidal}{comprimento de onda} enquanto
viaja --- é a dispersão.
\end{solution}

\begin{solution}{exo:b1:wave-propagation:11}
O modo $n$ tem forma $\sin(n\pi x/L)$. Um nó imposto em $x = L/3$ exige
$\sin(n\pi/3) = 0$, isto é,\ $n$ múltiplo de $3$: frequências $3f_1,
6f_1, \dots$ Dedilhar em $L/2$ dá à corda um deslocamento máximo
ali, para o qual o modo $n$ só pode contribuir se $\sin(n\pi/2) \neq 0$:
os modos pares, que têm um nó em $L/2$, nada recebem.
\end{solution}

\begin{solution}{exo:b1:wave-propagation:12}
$\lambda = c/f = \qty{3.0}{m}$. Diferença de \omterm{def:b1:wave-propagation:signal}{fase} $2\pi(2d)/\lambda + \pi$;
destrutiva quando ela vale $(2m + 1)\pi$: $2d/\lambda = m$,
$d = m\lambda/2 = 0, 1.5, 3.0, \dots$ m; construtiva quando $2d/\lambda =
m + \tfrac12$: $d = 0.75, 2.25, \dots$ m. As paredes e o teto do túnel refletem
a onda: um padrão estacionário com pontos mortos a cada meio \omterm{prop:b1:wave-propagation:sinusoidal}{comprimento de onda} ao longo
do percurso do carro.
\end{solution}

\begin{solution}{pb:b1:wave-propagation:1}
\textbf{1.} $[F/\mu] = M L T^{-2}/(M L^{-1}) = L^2 T^{-2}$: uma velocidade
ao quadrado.

\textbf{2.} Uma onda que vai para $-x$ depende de $t + x/c$, isto é, de
$\omega t + kx$; $k = \omega/c = 2\pi f/c$.

\textbf{3.} $c = \sqrt{0.90/1.0\times10^{-3}} = \qty{30}{m/s}$;
$\lambda = c/f = \qty{0.60}{m}$; $k = 2\pi/\lambda = \qty{10.5}{rad/m}$.

\textbf{4.} $\Delta\varphi = kd = 10.5 \times 0.15 = \qty{1.6}{rad} = \pi/2$.

\textbf{5.} $v_{\max} = A\omega = 2.0\times10^{-3} \times 2\pi \times 50 =
\qty{0.63}{m/s}$, cinquenta vezes menor que $c$: os pontos da corda se movem
devagar enquanto a forma corre.

\textbf{6.} Ela viaja para $+x$: \omterm{def:b1:wave-propagation:signal}{fase} $\omega t - kx$. Na extremidade
fixa o deslocamento total se anula em todo instante: $s_i(0, t) + s_r(0, t)
= A\cos\omega t + (-A)\cos\omega t = 0$.

\textbf{7.} $\cos(\omega t + kx) - \cos(\omega t - kx) =
-2\sin(\omega t)\sin(kx)$, logo $s = -2A\sin(kx)\sin(\omega t)$.

\textbf{8.} Nós: $\sin kx = 0$, $x = n\lambda/2$; ventres $x = (n +
\tfrac12)\lambda/2$; nós vizinhos separados por $\lambda/2$.

\textbf{9.} $\sin kL = 0$: $L = n\lambda/2$, isto é,\ $f_n = nc/2L$.

\textbf{10.} $n = 2L/\lambda = 2.4/0.60 = 4$: quatro fusos, um inteiro
--- ressonância.

\textbf{11.} $3$ fusos: $\lambda = 2L/3 = \qty{0.80}{m}$, $c = \lambda f
= \qty{40}{m/s}$, $F = \mu c^2 = \qty{1.6}{N}$. $6$ fusos: $\lambda =
\qty{0.40}{m}$, $c = \qty{20}{m/s}$, $F = \qty{0.40}{N}$.

\textbf{12.} $c = 2Lf/n = 2 \times 1.20 \times 50/3 = \qty{40}{m/s}$,
$\mu = F/c^2 = 1.6/1600 = \qty{1.0e-3}{kg/m} = \qty{1.0}{g/m}$.

\textbf{13.} $\mu = mgn^2/(4L^2f^2)$: $u(\mu)/\mu = \sqrt{(u_m/m)^2 +
(2u_L/L)^2} = \sqrt{0.61^2 + 0.33^2}\,\% = 0.7\%$; a massa é o que limita.

\textbf{14.} A fundamental. $c = 2Lf_1$: mi grave $2 \times 0.650 \times
82.4 = \qty{107}{m/s}$; mi agudo \qty{428}{m/s}.

\textbf{15.} $F = \mu c^2$: mi grave $6.0\times10^{-3} \times 107^2 =
\qty{69}{N}$; mi agudo $4.0\times10^{-4} \times 428^2 = \qty{73}{N}$ ---
tensões comparáveis. Baixar $f$ por um fator $4$ com $L$ e $\mu$ fixos
exigiria $16$ vezes menos tensão: uma corda mole e trepidante. Encapar
aumenta $\mu$ em vez disso.

\textbf{16.} $f_n = 2^{n/12}f_1$ e $f \propto 1/L$: $L_n = 2^{-n/12}L$,
distância da pestana $L - L_n = L(1 - 2^{-n/12})$: traste 1 \qty{36.5}{mm},
traste 5 \qty{163}{mm}, traste 7 \qty{216}{mm}, traste 12 \qty{325}{mm}.

\textbf{17.} Cada espaçamento vale $2^{-1/12} = 0.944$ vezes o anterior.

\textbf{18.} Fator $2^{3/12} = 1.19$ (três semitons); comprimento
$650 \times 2^{-3/12} = \qty{547}{mm}$.

\textbf{19.} Tocando em $L/2$: só modos pares, altura $2f_1 =
\qty{165}{Hz}$; em $L/3$: múltiplos de $3$, altura $3f_1 = \qty{247}{Hz}$.

\textbf{20.} Perto de uma extremidade, $\sin(n\pi x/L) \approx n\pi x/L$ cresce com
$n$: os modos altos são relativamente mais excitados --- som brilhante. Sobre a
boca (perto do meio) os modos pares são fracos e a fundamental domina
--- som doce.

\textbf{21.} $3f_1 = \qty{330}{Hz}$ bate a \qty{2}{Hz} com $2f_2$:
$f_2 = 166$ ou \qty{164}{Hz}, razão $1.509$ ou $1.491$ (a quinta
temperada vale $2^{7/12} = 1.498$).

\textbf{22.} O quarto harmônico, $4f$. Três batimentos: $4f = 437$ ou
\qty{443}{Hz}, $f = 109.25$ ou \qty{110.75}{Hz}.

\textbf{23.} Os batimentos aceleram ao apertar, logo a corda estava alta:
\qty{110.75}{Hz}. Ela deve baixar $0.75/110.75 = 0.68\%$, isto é,\ a
tensão deve cair $1.4\%$.

\textbf{24.} $f = (n/2L)\sqrt{F/\mu}$, logo $\ln f = \tfrac12\ln F + \text{const}$
e $\dd f/f = \tfrac12\,\dd F/F$.

\textbf{25.} $f_1 = \dfrac{1}{2L}\sqrt{\dfrac{F}{\mu}}$: o comprimento fixa
$\lambda_1 = 2L$, a tensão e a massa fixam $c$, os trastes escalam $L$ por
$2^{-n/12}$, e os batimentos revelam $f - f_{\text{ref}}$. Por trás de tudo, uma
relação: $\lambda f = c$, com a condição de \omterm{thm:b1:wave-propagation:standing}{onda estacionária} $L = n\lambda/2$.
\end{solution}
